\section{The \shadow Benchmark}
\label{sec:benchmark}

\subsection{Dataset Design and City Selection}
\label{sec:design}

\shadow contains $4{,}500$ RGB patches of size $384{\times}384$ pixels, sampled from National Agriculture Imagery Program (NAIP) orthoimagery across three U.S.\ cities: Chicago, Miami, and Phoenix. Each city contributes $1{,}500$ patches, split evenly across two native NAIP resolutions, $0.3$\,m and $0.6$\,m ground sampling distance ($750$ patches per city$\times$resolution pair). The three cities span distinct urban morphologies and climatic regimes, so the leave-one-city-out protocol in Section~\ref{sec:loco} measures transfer across genuine appearance shifts rather than nearby variants of the same scene type. Other potential confounders (sensor type, image format, file conventions) are held fixed by the shared NAIP source. We retain the two native resolutions rather than synthetically downsampling $0.3$\,m imagery, since NAIP releases them as separate native acquisitions with different sensor characteristics; collapsing these differences would understate the gap deployed systems face. Preserving both resolutions also supports future study of resolution shift under the same protocol. Within each city, patches are sampled uniformly inside the administrative boundary, with at most one patch per source NAIP tile to limit spatial autocorrelation. For each city-resolution pair, the $750$ patches are split into $450$ training, $150$ validation, and $150$ test images, fixed at the benchmark level and released with the dataset. To support resolution-paired analyses, we additionally co-locate $150$ coordinates per city imaged at both resolutions ($300$ paired patches per city). Sampling code, random seeds, and per-patch metadata (city, resolution, longitude/latitude, source tile ID) are released with the benchmark; within-city split distributional diagnostics are in Appendix~\ref{app:splits}.

\subsection{Annotation Pipeline}
\label{sec:annotation}

Pixel-level shadow masks are produced by $7$ trained annotators in a three-phase pipeline that jointly solves throughput and label quality. Phase~1 establishes the label distribution from scratch through human annotation. Phase~2 trains an interim shadow detector on Phase~1 labels and uses its predictions to seed subsequent batches, which annotators edit rather than redraw. Phase~3 is a retroactive quality-control pass that intersects each prior segment with the final-model prediction and returns disagreements for adjudication. Across all phases, every $100$-image batch includes $5$ gold images with frozen reference masks; Cohen's $\kappa$ on these images monitors annotator drift across batches. Each shadow segment carries one of six type labels (taxonomy in Appendix~\ref{app:data_stats}); the six-way typing is released for downstream analyses, but the LOCO benchmark in Section~\ref{sec:loco} evaluates binary masks (shadow vs.\ background). Tooling, per-phase throughput, inter-annotator agreement statistics, and Phase~3 adjudication rules are in Appendix~\ref{app:annotation}.

\begin{table}[t]
\centering
\small
\setlength{\tabcolsep}{2.4pt}
\caption{
\textbf{LOCO baseline performance across architectures and adaptation methods.} Two-class boundary-tolerant mIoU ($\pm 2$\,px). For each architecture-method pair we report mIoU pooled across the $6$ test cells per pair ($3$ held-out cities $\times$ $2$ native resolutions, $150$ paired test images per cell). \textit{Upper} is the within-city upper bound (per-architecture, pooled across the same cells). $\Delta$ is the mean gap to the upper bound (mIoU points; negative is degradation). \textit{Sig.} is the count of cells, out of $6$, with significant degradation under paired bootstrap ($p < 0.05$, $B = 10{,}000$). $^\dagger$High-resolution cells only ($3$ held-out cities); Sig.\ is given out of the evaluated cells. The \textit{Pooled} block pools MAMNet, OGLANet and DINOv3.
}
\label{tab:loco-main}
\begin{tabular}{l|cccc|cccc|cccc|cc}
\toprule
& \multicolumn{4}{c|}{\textit{LOCO mIoU}}
& \multicolumn{4}{c|}{$\Delta$ \textit{(mIoU)}}
& \multicolumn{4}{c|}{\textit{Sig.}}
& \multicolumn{2}{c}{\textit{Pooled}} \\
\textbf{Method} & \textbf{D} & \textbf{M} & \textbf{O} & \textbf{C}$^\dagger$
                & \textbf{D} & \textbf{M} & \textbf{O} & \textbf{C}$^\dagger$
                & \textbf{D} & \textbf{M} & \textbf{O} & \textbf{C}$^\dagger$
                & $\Delta$ & \textit{Sig./18} \\
\midrule
\textit{Upper} & $\mathit{77.6}$ & $\mathit{81.7}$ & $\mathit{74.9}$ & $\mathit{78.8}$ & --- & --- & --- & --- & --- & --- & --- & --- & --- & --- \\
\midrule
Vanilla   & 70.0 & 77.1 & 67.0 & 75.3 & $-7.6$ & $-4.6$ & $-7.9$ & $-3.5$ & 6 & 6 & 6 & 3/3 & $-6.68$ & 18 \\
FDA       & 69.4 & 77.0 & 67.0 & 75.1 & $-8.1$ & $-4.7$ & $-8.0$ & $-3.7$ & 6 & 6 & 6 & 3/3 & $-6.93$ & 18 \\
SegDeSiC  & 73.0 & 76.1 & 65.0 & -- & $-4.6$ & $-5.6$ & $-9.9$ & -- & 5 & 6 & 6 & -- & $-6.68$ & 17 \\
IIM       & 71.4 & 71.9 & 65.5 & -- & $-6.1$ & $-9.7$ & $-9.4$ & -- & 6 & 6 & 6 & -- & $-8.41$ & 18 \\
MRFP+     & 72.7 & 71.4 & 60.9 & -- & $-4.9$ & $-10.3$ & $-14.0$ & -- & 6 & 6 & 6 & -- & $-9.74$ & 18 \\
FADA      & 67.0 & 72.3 & 63.8 & -- & $-10.6$ & $-9.4$ & $-11.1$ & -- & 6 & 6 & 6 & -- & $-10.36$ & 18 \\
ISW       & 59.0 & 68.3 & 63.9 & -- & $-18.5$ & $-13.4$ & $-11.1$ & -- & 6 & 6 & 6 & -- & $-14.32$ & 18 \\
SAR$^\dagger$ & 68.4 & 73.1 & 65.4 & 75.3 & $-6.3$ & $-6.6$ & $-7.6$ & $-3.5$ & 3/3 & 3/3 & 3/3 & 3/3 & $-6.86$ & 9/9 \\
BSM$^\dagger$ & 66.2 & 72.8 & 62.3 & 75.7 & $-5.7$ & $-6.9$ & $-10.7$ & $-3.0$ & 2/2 & 3/3 & 3/3 & 3/3 & $-8.05$ & 8/8 \\
\midrule
\textit{Avg.\ $\Delta$} & --- & --- & --- & --- & $-8.6$ & $-8.2$ & $-10.2$ & --- & --- & --- & --- & --- & --- & --- \\
\textit{\# Sig./$42$} & --- & --- & --- & --- & --- & --- & --- & --- & 41 & 42 & 42 & --- & --- & --- \\
\bottomrule
\end{tabular}
\\[2pt]
\footnotesize
\textbf{D}: DINOv3 \quad \textbf{M}: MAMNet \quad \textbf{O}: OGLANet \quad \textbf{C}: CrossEarth. Avg.\ $\Delta$ and \# Sig.\ cover the first seven methods. --: not run; FADA, IIM, MRFP+, SegDeSiC and ISW are defined on a per-pixel head, which CrossEarth's query-based Mask2Former decoder lacks. CrossEarth \sib{} ($-4.4$, 3/3) is in Table~\ref{tab:loco-full-hr}.
\end{table}

\subsection{The LOCO Protocol}
\label{sec:loco}

We define the leave-one-city-out (LOCO) protocol over the city set $\mathcal{C} = \{\text{Chicago, Miami, Phoenix}\}$. For each holdout city $c^* \in \mathcal{C}$, the training set is sampled from the two non-holdout cities at a fixed $50/50$ ratio, yielding $n_{\text{train}} = 450$ images ($225$ per training city). The validation set is $n_{\text{val}} = 300$ images drawn from the same two training cities ($150$ per city), and the test set is the held-out city's $150$-image test pool. The within-city upper bound for each (city, resolution) is a model trained on that pair's $450$ training images and evaluated on its $150$ test images, so the LOCO-vs-upper-bound gap isolates geographic shift from training-set scale. All three holdout rotations are evaluated; results aggregate over rotations and the two native resolutions, giving $6$ test cells per (architecture, method) pair. The primary metric is boundary-tolerant mIoU, the mean of shadow and background IoU with a $\pm 2$-pixel don't-care band, since shadow boundaries are inherently imprecise to annotate at sub-pixel resolution; full metric definition and excluded-mask handling are in Appendix~\ref{app:metric}. Per-cell significance uses paired bootstrap on per-image $\Delta$ values ($B = 10{,}000$ resamples, two-tailed $p$-value against the within-city upper bound).

\subsection{Baselines}
\label{sec:baselines}

We evaluate four architectures spanning representative design families: \textbf{MAMNet}~\cite{mamnet} (CNN with criss-cross attention), \textbf{OGLANet}~\cite{oglanet} (CNN with multi-scale prediction-stage output aggregation), \textbf{DINOv3}~\cite{dinov3} (ViT with self-supervised pretraining, fine-tuned end-to-end), and \textbf{CrossEarth}~\cite{gong2025crossearth} (remote-sensing foundation model with a frozen DINOv2 backbone and a Mask2Former decoder). MAMNet, OGLANet and DINOv3 are each paired with nine configurations: \textbf{Vanilla} (source-only, no domain adaptation), \textbf{FDA}~\cite{fda}, \textbf{SegDeSiC}~\cite{segdesicnet}, \textbf{IIM}~\cite{yola}, \textbf{ISW}~\cite{isw}, \textbf{MRFP+}~\cite{mrfp}, \textbf{FADA}~\cite{fada}, \textbf{SAR}~\cite{niu2023sar}, and \textbf{BSM}~\cite{luo2023bsm}. Adaptation methods leave the architecture unchanged. The first seven are tested on all $7 \times 3 \times 3 \times 2 = 126$ (method, architecture, holdout city, resolution) configurations, and SAR, BSM and CrossEarth on the high-resolution cells, for $158$ LOCO test cells of $150$ images each. Architecture details, pretrained weight sources, and adaptation-method hyperparameters are in Appendix~\ref{app:training}. Table~\ref{tab:loco-main} summarizes results pooled across the three holdouts and two resolutions; per-cell numbers and confidence intervals are in Appendix~\ref{app:loco_extended}. $157$ of the $158$ test cells degrade significantly under geographic transfer ($p < 0.05$, paired bootstrap; $151$ after Bonferroni correction over the $158$-test family), with mean LOCO degradation of $-8.34$ mIoU and a worst single cell of $-28.50$ mIoU. A data-scaling test indicates that the gap is not a small-data effect: quadrupling the LOCO training set from $225$ to $900$ images narrows it by $0.80$ mIoU ($95\%$ CI $[0.46, 1.14]$), from $6.48$ to $5.50$, with the upper bound significantly ahead in every cell but one (high resolution, three seeds; Appendix~\ref{app:ladder}). Two findings in Table~\ref{tab:loco-main} shape the rest of the paper. First, the cost of an adaptation method depends on the architecture: ISW costs DINOv3 $-18.5$ mIoU against $-11.1$ on OGLANet, while MRFP$^+$ costs OGLANet $-14.0$ against $-4.9$ on DINOv3. Second, no adaptation method is uniformly safe: six of eight adaptation methods degrade more than Vanilla on average (five significantly, Holm $p < 0.05$), and every method except SegDeSiC degrades every cell it is evaluated on, on all four architectures. Without access to the target city, Table~\ref{tab:loco-main} gives a practitioner no basis for choosing an adaptation method. Section~\ref{sec:evaluation} asks \textit{where} in the architecture the failure lives and \textit{why} uniform corrections fail non-uniformly across them.

